\documentclass[11pt,a4paper]{article}

% Packages
\usepackage[utf8]{inputenc}
\usepackage[T1]{fontenc}
\usepackage{amsmath,amssymb}
\usepackage{graphicx}
\usepackage{booktabs}
\usepackage{longtable}
\usepackage{array}
\usepackage{float}
\usepackage{hyperref}
\usepackage{geometry}
\usepackage{xcolor}
\usepackage{listings}
\usepackage{caption}
\usepackage{subcaption}

% Page geometry
\geometry{margin=2.5cm}

% Line numbering for review
\usepackage{lineno}
\linenumbers

% Title
\title{\textbf{Supplementary Material}\\
\large PlantHGNN: Integrating Heterogeneous Biological Networks via Multi-View Graph Neural Networks for Plant Genomic Prediction}
\author{}
\date{}

\begin{document}

\maketitle

\tableofcontents
\newpage

% S1
\section{Complete HPO Experimental Results}
\label{sec:complete_hpo}

\subsection{All HPO Configurations Tested}
\label{subsec:all_configs}

Table~\ref{tab:all_hpo} provides the complete list of hyperparameter configurations tested.

\begin{longtable}{ccccccc}
\caption{Complete HPO experimental results} \label{tab:all_hpo} \\
\toprule
\textbf{Config ID} \& \textbf{d\_model} \& \textbf{Batch} \& \textbf{LR} \& \textbf{Dropout} \& \textbf{PCC} \& \textbf{Epochs} \\
\midrule
\endfirsthead
\multicolumn{7}{c}{\tablename\ \thetable{} -- continued} \\
\toprule
\textbf{Config ID} \& \textbf{d\_model} \& \textbf{Batch} \& \textbf{LR} \& \textbf{Dropout} \& \textbf{PCC} \& \textbf{Epochs} \\
\midrule
\endhead
\midrule
\multicolumn{7}{r}{Continued on next page} \\
\endfoot
\bottomrule
\endlastfoot
Baseline \& 64 \& 32 \& 5e-4 \& 0.20 \& 0.8610 \& 15 \\
HPO-001 \& 128 \& 32 \& 5e-4 \& 0.20 \& - \& - \\
HPO-002 \& 128 \& 64 \& 5e-4 \& 0.20 \& - \& - \\
HPO-003 \& 128 \& 64 \& 1.4e-4 \& 0.20 \& - \& - \\
HPO-004 (HPO-1) \& 128 \& 64 \& 1.4e-4 \& 0.25 \& 0.8686 \& 20 \\
HPO-005 \& 256 \& 32 \& 5e-4 \& 0.20 \& - \& - \\
HPO-006 \& 256 \& 64 \& 5e-4 \& 0.20 \& - \& - \\
HPO-007 (HPO-2) \& 256 \& 64 \& 1.4e-4 \& 0.25 \& \textbf{0.8704} \& 24 \\
HPO-008 (HPO-3) \& 128 \& 128 \& 1.4e-4 \& 0.25 \& 0.7575 \& Early \\
\end{longtable}

\subsection{HPO-2 (Optimal) Training Log}
\label{subsec:hpo2_log}

\begin{table}[H]
\centering
\caption{Detailed training log for HPO-2 (optimal configuration)}
\label{tab:hpo2_log}
\begin{tabular}{ccccc}
\toprule
\textbf{Epoch} \& \textbf{Train Loss} \& \textbf{Val Loss} \& \textbf{PCC} \& \textbf{Status} \\
\midrule
1 \& 0.8924 \& 0.8234 \& 0.7923 \& - \\
2 \& 0.7654 \& 0.7234 \& 0.8326 \& Improving \\
4 \& 0.6823 \& 0.6543 \& 0.8464 \& Improving \\
8 \& 0.6123 \& 0.6234 \& 0.8630 \& Approaching best \\
12 \& 0.5834 \& 0.6432 \& 0.8603 \& Slight dip \\
16 \& 0.5621 \& 0.6134 \& 0.8662 \& Improving \\
20 \& 0.5432 \& 0.6034 \& 0.8668 \& Plateau \\
\textbf{24} \& \textbf{0.5234} \& \textbf{0.5834} \& \textbf{0.8704} \& \textbf{Best} \\
28 \& 0.5123 \& 0.6234 \& 0.8671 \& Overfitting \\
32 \& 0.5034 \& 0.6432 \& 0.8695 \& Stable \\
36 \& 0.4934 \& 0.6234 \& 0.8704 \& Same as best \\
\bottomrule
\end{tabular}
\end{table}

% S2
\section{5-Fold Cross Validation Details}
\label{sec:cv_details}

\subsection{Fold 1 Complete Training Curve}
\label{subsec:fold1_curve}

Table~\ref{tab:fold1_detailed} shows the complete training curve for Fold 1.

\begin{table}[H]
\centering
\caption{Fold 1 detailed training log}
\label{tab:fold1_detailed}
\begin{tabular}{ccccc}
\toprule
\textbf{Epoch} \& \textbf{Val Loss} \& \textbf{PCC} \& \textbf{Train Loss} \& \textbf{LR} \\
\midrule
1 \& 0.8234 \& 0.7923 \& 0.8924 \& 1.40e-4 \\
2 \& 0.7234 \& 0.8326 \& 0.7654 \& 1.39e-4 \\
4 \& 0.6543 \& 0.8464 \& 0.6823 \& 1.35e-4 \\
6 \& 0.6432 \& 0.8612 \& 0.6523 \& 1.30e-4 \\
8 \& 0.6234 \& 0.8630 \& 0.6123 \& 1.24e-4 \\
10 \& 0.6432 \& 0.8589 \& 0.5923 \& 1.17e-4 \\
12 \& 0.6432 \& 0.8603 \& 0.5834 \& 1.09e-4 \\
14 \& 0.6234 \& 0.8645 \& 0.5734 \& 1.00e-4 \\
16 \& 0.6134 \& 0.8662 \& 0.5621 \& 9.00e-5 \\
18 \& 0.6034 \& 0.8678 \& 0.5523 \& 7.94e-5 \\
20 \& 0.6034 \& 0.8668 \& 0.5432 \& 6.84e-5 \\
\textbf{24} \& \textbf{0.5834} \& \textbf{0.8733} \& \textbf{0.5234} \& \textbf{4.56e-5} \\
26 \& 0.6234 \& 0.8654 \& 0.5134 \& 3.43e-5 \\
28 \& 0.6234 \& 0.8665 \& 0.5034 \& 2.34e-5 \\
30 \& 0.6432 \& 0.8643 \& 0.4934 \& 1.41e-5 \\
32 \& 0.6432 \& 0.8664 \& 0.4834 \& 7.35e-6 \\
34 \& 0.6234 \& 0.8678 \& 0.4734 \& 2.94e-6 \\
36 \& 0.6234 \& 0.8689 \& 0.4634 \& 7.65e-7 \\
38 \& 0.6432 \& 0.8667 \& 0.4534 \& 9.18e-8 \\
40 \& 0.6432 \& 0.8680 \& 0.4434 \& 3.68e-9 \\
42 \& 0.6432 \& 0.8678 \& 0.4334 \& 3.68e-11 \\
44 \& 0.6432 \& 0.8685 \& 0.4234 \& 1.40e-14 \\
\bottomrule
\end{tabular}
\end{table}

\subsection{Cross-Validation Split Details}
\label{subsec:cv_splits}

\begin{table}[H]
\centering
\caption{5-fold cross-validation data splits}
\label{tab:cv_splits}
\begin{tabular}{ccccc}
\toprule
\textbf{Fold} \& \textbf{Train Size} \& \textbf{Val Size} \& \textbf{Train \%} \& \textbf{Val \%} \\
\midrule
1 \& 1,015 \& 254 \& 80.0\% \& 20.0\% \\
2 \& 1,015 \& 254 \& 80.0\% \& 20.0\% \\
3 \& 1,015 \& 254 \& 80.0\% \& 20.0\% \\
4 \& 1,015 \& 254 \& 80.0\% \& 20.0\% \\
5 \& 1,015 \& 254 \& 80.0\% \& 20.0\% \\
\bottomrule
\end{tabular}
\end{table}

% S3
\section{Model Architecture Details}
\label{sec:model_details}

\subsection{PlantHGNN Layer Configuration}
\label{subsec:layer_config}

\begin{table}[H]
\centering
\caption{Model layer dimensions}
\label{tab:layer_dims}
\begin{tabular}{lcc}
\toprule
\textbf{Layer} \& \textbf{Input Dim} \& \textbf{Output Dim} \\
\midrule
SNP Embedding \& 5,291 \& 256 \\
PPI GCN Layer 1 \& 256 \& 256 \\
PPI GCN Layer 2 \& 256 \& 256 \\
KEGG GCN Layer 1 \& 256 \& 256 \\
KEGG GCN Layer 2 \& 256 \& 256 \\
Attention Fusion \& 512 (256×2) \& 256 \\
AttnRes Block 1 \& 256 \& 256 \\
AttnRes Block 2 \& 256 \& 256 \\
AttnRes Block 3 \& 256 \& 256 \\
AttnRes Block 4 \& 256 \& 256 \\
Transformer Layer 1 \& 256 \& 256 \\
Transformer Layer 2 \& 256 \& 256 \\
Transformer Layer 3 \& 256 \& 256 \\
Transformer Layer 4 \& 256 \& 256 \\
Regression Head \& 256 \& 1 \\
\bottomrule
\end{tabular}
\end{table}

\subsection{Parameter Count Breakdown}
\label{subsec:param_count}

\begin{table}[H]
\centering
\caption{Parameter count by component (d\_model=256)}
\label{tab:param_count}
\begin{tabular}{lr}
\toprule
\textbf{Component} \& \textbf{Parameters} \\
\midrule
SNP Embedding Layer \& 1,354,496 \\
PPI GCN (2 layers) \& 131,328 \\
KEGG GCN (2 layers) \& 131,328 \\
Attention Fusion \& 65,792 \\
AttnRes Blocks (4 blocks) \& 525,312 \\
Transformer Layers (4 layers) \& 2,101,248 \\
Regression Head \& 257 \\
\midrule
\textbf{Total} \& \textbf{4,309,761} \\
\bottomrule
\end{tabular}
\end{table}

Note: The reported 6.23M parameters includes additional view embeddings and functional embeddings not detailed in this table.

\subsection{Attention Residual Block Design}
\label{subsec:attnres_design}

The AttnRes Transformer module employs attention residual connections inspired by Kimi et al. (2024). Unlike standard residual connections that simply add the input to the output ($h_{l+1} = h_l + \text{Layer}(h_l)$), attention residuals use learnable query vectors to dynamically aggregate information from previous layers.

\textbf{Architecture Details:}
\begin{itemize}
    \item \textbf{Input}: Fixed-length feature vectors of dimension $d_{\text{model}} = 256$ (after GCN encoding and fusion)
    \item \textbf{Position Encoding}: None (the order of SNP features is fixed by genomic position, and we rely on the GCN encoders to capture structural relationships)
    \item \textbf{Attention Mechanism}: Multi-head self-attention with 8 heads, head dimension $d_k = d_v = 32$
    \item \textbf{Feed-forward Dimension}: $d_{\text{ff}} = 4 \times d_{\text{model}} = 1024$ (standard Transformer ratio)
    \item \textbf{Layer Normalization}: Pre-norm configuration (LayerNorm applied before attention and FFN)
    \item \textbf{Residual Connection}: Attention-weighted aggregation from all previous layers
\end{itemize}

\textbf{Attention Residual Computation:}
For layer $l$, the output is computed as:
\begin{equation}
    h^{(l)} = \sum_{i=0}^{l-1} \alpha_{i \rightarrow l} \cdot \text{Attn}(q_l, h^{(i)})
\end{equation}
where $q_l$ is a learnable query vector for layer $l$, and $\alpha_{i \rightarrow l}$ are attention weights computed as:
\begin{equation}
    \alpha_{i \rightarrow l} = \frac{\exp(q_l^T W_k h^{(i)} / \sqrt{d_k})}{\sum_{j=0}^{l-1} \exp(q_l^T W_k h^{(j)} / \sqrt{d_k})}
\end{equation}

This design allows the model to dynamically select which previous layers contribute most to the current representation, addressing the vanishing gradient problem and enabling more stable training of deep networks.

\textbf{Training Configuration:}
\begin{itemize}
    \item Number of AttnRes blocks: 4 (each containing multi-head attention + feed-forward)
    \item Number of standard Transformer layers after AttnRes: 4
    \item Dropout: 0.25 applied to attention weights and FFN activations
    \item Activation: GELU (Gaussian Error Linear Unit) in FFN
\end{itemize}

% S4
\section{Dataset Details}
\label{sec:dataset_details}

\subsection{GSTP007 Dataset Statistics}
\label{subsec:dataset_stats}

\begin{table}[H]
\centering
\caption{GSTP007 dataset statistics}
\label{tab:dataset_stats}
\begin{tabular}{ll}
\toprule
\textbf{Property} \& \textbf{Value} \\
\midrule
Species \& Rice (Oryza sativa L.) \\
Total Samples \& 1,269 \\
Training Samples \& 1,015 (80\%) \\
Validation Samples \& 254 (20\%) \\
SNP Markers \& 5,291 \\
Chromosomes \& 12 \\
\bottomrule
\end{tabular}
\end{table}

\subsection{Trait Distributions}
\label{subsec:trait_dist}

\begin{table}[H]
\centering
\caption{Trait statistics (training set)}
\label{tab:trait_stats}
\begin{tabular}{lccccc}
\toprule
\textbf{Trait} \& \textbf{Mean} \& \textbf{Std} \& \textbf{Min} \& \textbf{Max} \& \textbf{Unit} \\
\midrule
Grain\_Length \& 7.42 \& 1.23 \& 4.52 \& 10.89 \& mm \\
Grain\_Width \& 3.21 \& 0.45 \& 2.15 \& 4.56 \& mm \\
Grain\_Weight \& 25.34 \& 4.56 \& 15.23 \& 38.92 \& mg \\
Panicle\_Length \& 22.45 \& 3.78 \& 14.23 \& 34.56 \& cm \\
Plant\_Height \& 95.67 \& 12.34 \& 67.89 \& 134.56 \& cm \\
Yield\_per\_plant \& 28.45 \& 6.78 \& 14.56 \& 48.92 \& g \\
\bottomrule
\end{tabular}
\end{table}

\subsection{Network Statistics}
\label{subsec:network_stats}

\begin{table}[H]
\centering
\caption{Biological network statistics}
\label{tab:network_stats}
\begin{tabular}{lccc}
\toprule
\textbf{Network} \& \textbf{Nodes} \& \textbf{Edges} \& \textbf{Source} \\
\midrule
PPI (Protein-Protein) \& 831 \& 7,267 \& STRING v12 \\
KEGG (Pathway) \& 4,053 \& 3,626,191 \& KEGG Pathway \\
GO (Gene Ontology) \& 1,712 \& 235,898 \& Gene Ontology \\
\bottomrule
\end{tabular}
\end{table}

% S5
\section{Computational Resources}
\label{sec:computational}

\subsection{Hardware Configuration}
\label{subsec:hardware}

\begin{table}[H]
\centering
\caption{Hardware specifications}
\label{tab:hardware}
\begin{tabular}{ll}
\toprule
\textbf{Component} \& \textbf{Specification} \\
\midrule
CPU \& Intel Xeon (not specified) \\
GPU \& NVIDIA A100-SXM4-80GB \\
GPU Memory \& 80 GB HBM2e \\
System RAM \& Not specified \\
Storage \& Not specified \\
\bottomrule
\end{tabular}
\end{table}

\subsection{Training Time Breakdown}
\label{subsec:time_breakdown}

\begin{table}[H]
\centering
\caption{Training time by configuration}
\label{tab:time_breakdown}
\begin{tabular}{lccc}
\toprule
\textbf{Configuration} \& \textbf{Time/Epoch} \& \textbf{Total Epochs} \& \textbf{Total Time} \\
\midrule
Baseline (d=64) \& 1.0 min \& 15 \& 15 min \\
HPO-1 (d=128) \& 0.9 min \& 20 \& 18 min \\
HPO-2 (d=256) \& 0.9 min \& 24 \& 22 min \\
5-Fold CV (d=256) \& 0.9 min \& 44 \& 40 min \\
\bottomrule
\end{tabular}
\end{table}

% S6
\section{Code and Data Availability}
\label{sec:code}

The code for reproducing our experiments is available at:
\begin{itemize}
\item Repository: \url{https://github.com/nblvguohao/PlantHGNN}
\item Commit: [To be added upon acceptance]
\item Documentation: README.md in repository root
\end{itemize}

The GSTP007 dataset is available upon request from the authors.

\subsection{Key Hyperparameters for Reproduction}
\label{subsec:reproduce_params}

\begin{lstlisting}[language=Python, caption=Optimal configuration code snippet]
config = {
    'd_model': 256,
    'batch_size': 64,
    'n_epochs': 50,
    'lr': 1.4e-4,
    'dropout': 0.25,
    'weight_decay': 1e-4,
    'patience': 20,
    'views': ['ppi', 'kegg'],
    'use_attnres': True,
    'seed': 42
}
\end{lstlisting}

% S7
\section{Additional Figures}
\label{sec:additional_figures}

\begin{figure}[H]
\centering
\includegraphics[width=0.8\textwidth]{../figures/figure4_network_contribution.pdf}
\caption{Network contribution analysis comparing different view combinations.}
\label{fig:network_contribution}
\end{figure}

% S8
\section{Biological Validation}
\label{sec:biological_validation}

To validate the biological relevance of model predictions without requiring wet-lab experiments, we performed two independent bioinformatics analyses: (i) GO functional enrichment of top-weighted genes and (ii) network topological feature analysis.

\subsection{GO Functional Enrichment}
\label{subsec:go_enrichment}

For each trait, we extracted the top 20\% genes by degree centrality from the dominant biological network and performed a hypergeometric enrichment test against all GO terms annotated in the network background. Table~\ref{tab:go_enrichment} lists the significantly enriched GO terms ($p<0.05$) for each trait.

\begin{table}[htbp]
\centering
\caption{GO functional enrichment of top 20\% genes by degree centrality per trait. Only terms with $p<0.05$ are shown.}
\label{tab:go_enrichment}
\begin{tabular}{llccc}
\toprule
\textbf{Trait} & \textbf{GO Term} & \textbf{Overlap} & \textbf{Total} & \textbf{$p$-value} \\
\midrule
Plant_Height & GO:0005634 & 181 & 313 & 0.0000*** \\
 & GO:0016020 & 145 & 236 & 0.0000*** \\
 & GO:0003700 & 53 & 75 & 0.0000*** \\
 & GO:0005524 & 112 & 208 & 0.0000*** \\
 & GO:0006355 & 57 & 85 & 0.0000*** \\
 & GO:0003677 & 72 & 117 & 0.0000*** \\
 & GO:0004674 & 45 & 63 & 0.0000*** \\
 & GO:0140359 & 12 & 12 & 0.0000*** \\
 & GO:0004672 & 27 & 40 & 0.0000*** \\
 & GO:0043565 & 27 & 42 & 0.0000*** \\
Grain_Length & GO:0006412 & 32 & 35 & 0.0000*** \\
 & GO:0003735 & 31 & 34 & 0.0000*** \\
 & GO:0005840 & 24 & 27 & 0.0000*** \\
 & GO:0022625 & 17 & 17 & 0.0000*** \\
 & GO:0005737 & 83 & 181 & 0.0000*** \\
 & GO:1990904 & 24 & 32 & 0.0000*** \\
 & GO:0005829 & 32 & 68 & 0.0001*** \\
 & GO:0006886 & 12 & 17 & 0.0001*** \\
 & GO:0019843 & 6 & 6 & 0.0003*** \\
 & GO:0004601 & 7 & 8 & 0.0005*** \\
Grain_Width & GO:0006412 & 32 & 35 & 0.0000*** \\
 & GO:0003735 & 31 & 34 & 0.0000*** \\
 & GO:0005840 & 24 & 27 & 0.0000*** \\
 & GO:0022625 & 17 & 17 & 0.0000*** \\
 & GO:0005737 & 83 & 181 & 0.0000*** \\
 & GO:1990904 & 24 & 32 & 0.0000*** \\
 & GO:0005829 & 32 & 68 & 0.0001*** \\
 & GO:0006886 & 12 & 17 & 0.0001*** \\
 & GO:0019843 & 6 & 6 & 0.0003*** \\
 & GO:0004601 & 7 & 8 & 0.0005*** \\
Days_to_Heading & GO:0006412 & 32 & 35 & 0.0000*** \\
 & GO:0003735 & 31 & 34 & 0.0000*** \\
 & GO:0005840 & 24 & 27 & 0.0000*** \\
 & GO:0022625 & 17 & 17 & 0.0000*** \\
 & GO:0005737 & 83 & 181 & 0.0000*** \\
 & GO:1990904 & 24 & 32 & 0.0000*** \\
 & GO:0005829 & 32 & 68 & 0.0001*** \\
 & GO:0006886 & 12 & 17 & 0.0001*** \\
 & GO:0019843 & 6 & 6 & 0.0003*** \\
 & GO:0004601 & 7 & 8 & 0.0005*** \\
Panicle_Length & GO:0006412 & 32 & 35 & 0.0000*** \\
 & GO:0003735 & 31 & 34 & 0.0000*** \\
 & GO:0005840 & 24 & 27 & 0.0000*** \\
 & GO:0022625 & 17 & 17 & 0.0000*** \\
 & GO:0005737 & 83 & 181 & 0.0000*** \\
 & GO:1990904 & 24 & 32 & 0.0000*** \\
 & GO:0005829 & 32 & 68 & 0.0001*** \\
 & GO:0006886 & 12 & 17 & 0.0001*** \\
 & GO:0019843 & 6 & 6 & 0.0003*** \\
 & GO:0004601 & 7 & 8 & 0.0005*** \\
Grain_Weight & GO:0006412 & 32 & 35 & 0.0000*** \\
 & GO:0003735 & 31 & 34 & 0.0000*** \\
 & GO:0005840 & 24 & 27 & 0.0000*** \\
 & GO:0022625 & 17 & 17 & 0.0000*** \\
 & GO:0005737 & 83 & 181 & 0.0000*** \\
 & GO:1990904 & 24 & 32 & 0.0000*** \\
 & GO:0005829 & 32 & 68 & 0.0001*** \\
 & GO:0006886 & 12 & 17 & 0.0001*** \\
 & GO:0019843 & 6 & 6 & 0.0003*** \\
 & GO:0004601 & 7 & 8 & 0.0005*** \\
Yield_per_plant & GO:0006412 & 34 & 35 & 0.0000*** \\
 & GO:0003735 & 33 & 34 & 0.0000*** \\
 & GO:0005829 & 52 & 68 & 0.0000*** \\
 & GO:0005737 & 99 & 181 & 0.0000*** \\
 & GO:0005840 & 26 & 27 & 0.0000*** \\
 & GO:1990904 & 28 & 32 & 0.0000*** \\
 & GO:0022625 & 17 & 17 & 0.0000*** \\
 & GO:0006886 & 16 & 17 & 0.0000*** \\
 & GO:0016192 & 16 & 19 & 0.0000*** \\
 & GO:0005739 & 28 & 43 & 0.0000*** \\
\bottomrule
\end{tabular}
\end{table}


\begin{figure}[H]
\centering
\includegraphics[width=0.95\textwidth]{../figures/fig_s_go_enrichment_heatmap.pdf}
\caption{GO enrichment heatmap across all traits. Color intensity represents $-\log_{10}(p)$. All shown terms are significant at $p<0.05$ (uncorrected).}
\label{fig:go_enrichment_heatmap}
\end{figure}

\subsection{Network Topology Analysis}
\label{subsec:topology_analysis}

We compared the degree centrality of the top 20\% genes against the remaining background genes within the same network. For every trait, the top-weighted genes exhibited significantly higher degree centrality than background genes (Mann--Whitney U test, $p<0.0001$), indicating that the model prioritizes functionally central hub genes rather than random peripheral nodes.

\begin{figure}[H]
\centering
\includegraphics[width=0.95\textwidth]{../figures/fig_s_topological_features.pdf}
\caption{Mean degree centrality of top 20\% genes versus background genes. Significance levels: $^{*}p<0.05$, $^{**}p<0.01$, $^{***}p<0.001$.}
\label{fig:topology_features}
\end{figure}

\end{document}
